% !TeX root = main.tex
%
% Theory section for the Pavillon 3D Gaussian Splatting report.
% Compile with:  latexmk -pdf main.tex   (see main.tex in this directory)

\section{Theoretical Background}
\label{sec:theory}

This section develops the theory underlying the reconstruction pipeline: the
Gaussian-splatting scene representation and its differentiable rasteriser
(\S\ref{sec:3dgs}), the structure-from-motion stage that supplies camera poses
(\S\ref{sec:sfm}), the adaptive density control that grows the model
(\S\ref{sec:densification}), and the four regularisation mechanisms that adapt the
method to our capture: monocular depth priors, spatial bounding, floater
suppression, and per-image appearance modelling
(\S\ref{sec:sparse}--\S\ref{sec:appearance}). Throughout, we highlight
\emph{why} each component is required by the specific difficulty of our data: a
single-sided, low-parallax, close-range capture of a planar carved relief.

%---------------------------------------------------------------------------
\subsection{Scene representation: 3D Gaussian Splatting}
\label{sec:3dgs}

Neural radiance fields represent a scene as a continuous function evaluated by
ray marching \cite{mildenhall2020nerf}, which couples rendering cost to sample
count. 3D Gaussian Splatting (3DGS) \cite{kerbl2023gaussian} instead adopts an
explicit, unstructured set of anisotropic Gaussian primitives that are
\emph{rasterised}, giving real-time rendering while remaining fully
differentiable.

The scene is a set $\mathcal{G}=\{g_i\}_{i=1}^{N}$ where each primitive carries a
mean $\bm{\mu}_i\in\mathbb{R}^3$, a covariance $\bm{\Sigma}_i\in\mathbb{R}^{3\times3}$,
an opacity $o_i\in\mathbb{R}$ and view-dependent colour coefficients. The spatial
support is the unnormalised Gaussian
\begin{equation}
  G_i(\bm{x}) \;=\; \exp\!\Big(-\tfrac{1}{2}(\bm{x}-\bm{\mu}_i)^{\!\top}\,
  \bm{\Sigma}_i^{-1}\,(\bm{x}-\bm{\mu}_i)\Big).
  \label{eq:gaussian}
\end{equation}

\paragraph{Parameterisation of the covariance.}
A covariance matrix must remain symmetric positive semi-definite, which gradient
descent on its six free entries does not guarantee. 3DGS therefore factorises it
into a rotation and an axis-aligned scaling,
\begin{equation}
  \bm{\Sigma}_i \;=\; \bm{R}(\bm{q}_i)\,\bm{S}_i\,\bm{S}_i^{\!\top}\,\bm{R}(\bm{q}_i)^{\!\top},
  \qquad
  \bm{S}_i=\operatorname{diag}\!\big(\exp(\bm{s}_i)\big),
  \label{eq:covariance}
\end{equation}
optimising a unit quaternion $\bm{q}_i$ and \emph{log}-scales
$\bm{s}_i\in\mathbb{R}^3$. The exponential keeps scales strictly positive and makes
the parameterisation multiplicative, so a single learning rate behaves sensibly
across the wide range of primitive sizes in a scene. Opacity is likewise stored as
a logit and mapped through a sigmoid, $\alpha_i=\sigma(o_i)\in(0,1)$. These
reparameterisations matter operationally: our floater diagnostics
(\S\ref{sec:floaters}) threshold $\sigma(o_i)$ and $\max\exp(\bm{s}_i)$ rather
than the raw parameters.

\paragraph{Projection and the 2D covariance.}
Rendering requires the image-space footprint of each primitive. Under a viewing
transform $\bm{W}$ (world-to-camera) followed by a projective map, the projection
is nonlinear, so 3DGS adopts the EWA splatting approximation
\cite{zwicker2001ewa,zwicker2001surface}: linearise the projective map about
$\bm{\mu}_i$ with Jacobian $\bm{J}$ and propagate the covariance,
\begin{equation}
  \bm{\Sigma}'_i \;=\; \bm{J}\,\bm{W}\,\bm{\Sigma}_i\,\bm{W}^{\!\top}\bm{J}^{\!\top},
  \label{eq:proj}
\end{equation}
retaining the upper-left $2\times2$ block as the screen-space covariance. The
first-order approximation is accurate when a primitive is small relative to its
depth --- an assumption that our extreme close-up views stress, since a primitive
covering a large solid angle has a genuinely non-Gaussian footprint.

\paragraph{Compositing.}
Colour is accumulated by front-to-back alpha compositing, the discretisation of
the volume-rendering integral under an absorption-emission model
\cite{max1995optical}. For a pixel $\bm{p}$, with primitives sorted by depth,
\begin{equation}
  \bm{C}(\bm{p}) \;=\; \sum_{i=1}^{N} \bm{c}_i(\bm{d})\,\hat{\alpha}_i(\bm{p})
  \prod_{j<i}\big(1-\hat{\alpha}_j(\bm{p})\big),
  \qquad
  \hat{\alpha}_i(\bm{p}) = \alpha_i\,G'_i(\bm{p}),
  \label{eq:compositing}
\end{equation}
where $G'_i$ is the projected 2D Gaussian~\eqref{eq:proj} evaluated at the pixel
and $\bm{c}_i(\bm{d})$ is the colour emitted toward the viewing direction
$\bm{d}$, expanded in spherical harmonics,
$\bm{c}_i(\bm{d})=\sum_{\ell=0}^{L}\sum_{m=-\ell}^{\ell}\bm{k}_{i,\ell m}\,Y_{\ell m}(\bm{d})$.
We use $L=3$ (16 coefficients per colour channel) and, following common practice,
introduce the bands progressively during training so that low-frequency albedo is
fitted before view-dependent effects.

Because \eqref{eq:compositing} is a product of per-primitive terms, every quantity
is differentiable and gradients flow to positions, covariances, opacities and
colours simultaneously. This is also the source of a characteristic failure mode:
the same pixel colour can be explained either by a correctly placed opaque
primitive or by a stack of misplaced semi-transparent ones. Nothing in
\eqref{eq:compositing} prefers the former, which is precisely why explicit
regularisation (\S\ref{sec:floaters}) is necessary.

\paragraph{Photometric objective.}
Training minimises a combination of an $\ell_1$ term and a structural
dissimilarity term \cite{wang2004ssim},
\begin{equation}
  \mathcal{L}_{\text{photo}} \;=\; (1-\lambda)\,\big\lVert \hat{\bm{I}}-\bm{I}\big\rVert_1
  \;+\; \lambda\,\big(1-\mathrm{SSIM}(\hat{\bm{I}},\bm{I})\big),
  \label{eq:photo}
\end{equation}
with $\lambda=0.2$. The SSIM term supplies a spatially structured gradient that a
purely per-pixel loss lacks, which matters for high-frequency relief detail.
Optimisation uses Adam \cite{kingma2015adam} with per-parameter-group learning
rates; the positional learning rate is decayed exponentially so primitives settle
rather than jitter, a detail that proved decisive for sharpness in our runs.

%---------------------------------------------------------------------------
\subsection{Camera calibration: incremental versus global SfM}
\label{sec:sfm}

3DGS optimises geometry \emph{given} camera poses; those come from
structure-from-motion. The incremental pipeline of COLMAP
\cite{schoenberger2016sfm} extracts SIFT features \cite{lowe2004sift}, matches
them, then seeds a two-view reconstruction and repeatedly registers one image at a
time, triangulating and running bundle adjustment. Incremental growth is robust on
well-connected sets but is greedy: registration proceeds only where the current
model already provides sufficient $2$D--$3$D correspondences, so a weakly
connected region can be abandoned entirely.

Global SfM instead estimates all camera poses simultaneously from the epipolar
graph. GLOMAP \cite{pan2024glomap} performs global rotation averaging followed by
a joint global positioning step that estimates camera centres and points together,
avoiding the classical fragility of separate translation averaging. Because it
considers the whole view graph at once, it does not inherit the incremental
pipeline's order dependence.

This distinction was decisive for our capture. On the single-orbit Pavillon clip,
incremental COLMAP registered $82$ images whereas GLOMAP registered $181$ of
$193$ from identical features and matches --- a $2.2\times$ coverage gain. Low
image overlap produces exactly the weakly connected graph that penalises greedy
registration, and coverage is the binding constraint on downstream quality: a
surface observed by no registered camera cannot be reconstructed at all.

%---------------------------------------------------------------------------
\subsection{Adaptive density control}
\label{sec:densification}

A fixed primitive count cannot match a scene's spatially varying complexity, so
3DGS \emph{adapts} the set during optimisation \cite{kerbl2023gaussian}. The
driving signal is the magnitude of the loss gradient with respect to the
\emph{projected} (screen-space) mean, averaged over the views in which a primitive
was visible,
\begin{equation}
  \tau_i \;=\; \Big\lVert \frac{\partial \mathcal{L}}{\partial \bm{\mu}'_i}\Big\rVert.
\end{equation}
A large $\tau_i$ indicates that the renderer would reduce the loss by moving the
primitive --- a signature of under-reconstruction. Where $\tau_i$ exceeds a
threshold, the model is refined in one of two ways:
\begin{itemize}
  \item \textbf{Cloning} (primitive smaller than a scene-scale-relative bound):
        duplicate it and displace along the positional gradient, adding capacity
        to an under-populated region;
  \item \textbf{Splitting} (primitive larger than that bound): replace it by two
        smaller primitives sampled from its own distribution, with scales divided
        by a constant, refining an over-large blob.
\end{itemize}
Periodically all opacities are reset toward zero; primitives that are genuinely
supported by the data recover, while those that merely accumulated as
semi-transparent clutter fall below the pruning threshold and are removed.

Two scene-scale dependencies deserve emphasis, because both caused real failures
in this project. First, the positional learning rate is scaled by the scene
extent; second, the clone-versus-split decision compares primitive size against a
scene-scale-relative bound. Estimating the scene extent from the raw point cloud
makes it hostage to SfM outliers: a single stray point inflates the extent, which
simultaneously raises the positional learning rate (primitives jitter, producing
blur) and raises the split threshold (primitives are only ever cloned, never
shrink). We therefore define the extent from \emph{camera centres} rather than
points, following the original spatial-scale convention.

%---------------------------------------------------------------------------
\subsection{The sparse-view problem and monocular depth priors}
\label{sec:sparse}

Equation~\eqref{eq:compositing} is severely underdetermined when views are few or
share little parallax. A set of primitives can reproduce every training image
exactly while placing mass at physically wrong depths, because depth is only
constrained by \emph{disagreement} between views. With a single-sided capture the
disagreement is small, and the photometric loss alone cannot distinguish a correct
surface from a family of depth-shifted alternatives.

The established remedy is to import geometric knowledge from a monocular depth
network \cite{ranftl2022midas,yang2024depthanythingv2} and use it to regularise
the rendered depth, as in FSGS \cite{zhu2024fsgs}, SparseGS
\cite{xiong2024sparsegs} and DNGaussian \cite{li2024dngaussian}. The rendered
depth is obtained by compositing depth exactly as colour is composited in
\eqref{eq:compositing}:
\begin{equation}
  \hat{D}(\bm{p}) \;=\; \sum_i d_i\,\hat{\alpha}_i(\bm{p}) \prod_{j<i}\big(1-\hat{\alpha}_j(\bm{p})\big).
\end{equation}

\paragraph{Why the loss must be affine-invariant.}
Monocular networks trained on mixed datasets predict depth only \emph{up to an
unknown scale and shift} \cite{ranftl2022midas}; the prediction lives in a
disparity-like space with no metric meaning. Penalising
$\lVert \hat{D}-D_{\text{mono}}\rVert$ directly would therefore impose an
arbitrary global scale on the reconstruction and fight the SfM geometry. The loss
must be invariant to the ambiguity, a requirement dating to the scale-invariant
objective of Eigen et al.~\cite{eigen2014depth}.

We adopt the Pearson correlation \cite{pearson1895} between rendered and predicted
depth over the pixels of a view,
\begin{equation}
  \mathcal{L}_{\text{depth}} \;=\; 1-\rho\big(\hat{D},\,D_{\text{mono}}\big),
  \qquad
  \rho(X,Y)=\frac{\operatorname{Cov}(X,Y)}{\sigma_X\,\sigma_Y}.
  \label{eq:pearson}
\end{equation}
For any $a>0$ and $b\in\mathbb{R}$ we have $\rho(aX+b,Y)=\rho(X,Y)$, so
\eqref{eq:pearson} is exactly invariant to the affine ambiguity: it constrains the
depth \emph{ordering} of the scene while leaving scale and offset to be fixed by
the cameras. Note the sign convention matters --- since the network emits a
disparity-like quantity (large means \emph{near}), it must be inverted before use,
or the loss would drive the geometry to the reflection of the correct surface.

The prior is ramped in only after the photometric term has established coarse
structure, so that a noisy monocular estimate does not dominate early
optimisation.

%---------------------------------------------------------------------------
\subsection{Floaters: spatial bounding and opacity/scale control}
\label{sec:floaters}

``Floaters'' are primitives occupying empty space, typically faint and often
large, which reproduce training views by acting as a diffuse haze near the cameras
while corrupting novel views. As noted after \eqref{eq:compositing}, the
compositing model gives no intrinsic preference for compact opaque geometry over
diffuse semi-transparent clutter; in the sparse-view regime such clutter is an
easy local minimum. In NeRF-style methods the standard countermeasure is the
distortion loss of Mip-NeRF~360 \cite{barron2022mipnerf360}, which penalises
weight dispersed along a ray and thereby encourages compact ray termination;
Mip-Splatting \cite{yu2024mipsplatting} addresses the related scale-aliasing
artefacts that appear when primitives are observed at sampling rates far from
those seen in training.

We use two direct, geometrically motivated constraints. First, a \emph{scene
bound}: the capture took place in a room, so admissible geometry lies within the
axis-aligned bounding box of the camera centres and the SfM points, dilated by a
margin $m$,
\begin{equation}
  \mathcal{B} \;=\; \big[\bm{l}-m(\bm{h}-\bm{l}),\;\bm{h}+m(\bm{h}-\bm{l})\big],
  \qquad
  \bm{l}=\min(\mathcal{P}\cup\mathcal{C}),\;\;
  \bm{h}=\max(\mathcal{P}\cup\mathcal{C}).
\end{equation}
Primitives leaving $\mathcal{B}$ have their opacity driven to zero and are then
removed by the existing opacity prune. Second, an explicit \emph{scale and opacity
criterion}: primitives whose largest scale exceeds a fraction of the scene extent,
or whose opacity falls below a threshold, are suppressed and finally pruned.

Suppressing rather than deleting mid-training is deliberate: the density-control
state of \S\ref{sec:densification} is indexed per primitive, so removing entries
externally would desynchronise it. Driving opacity to $\sigma(-20)\approx 2\times10^{-9}$
changes no tensor shape, makes the primitive invisible, and lets the existing
prune reclaim it.

Empirically these constraints removed $\approx\!18\text{--}20\%$ of primitives.
Crucially, the diagnostic revealed which mechanism did the work: only $0.01\%$ of
baseline primitives lay outside $\mathcal{B}$, so the bound acted as a safety net,
whereas $18.2\%$ had opacity below $0.02$ and the largest scales spanned the whole
scene. The floaters here were in-volume haze rather than runaway geometry, and it
was the opacity/scale criterion that eliminated them --- median opacity rose from
$0.51$ to $0.88$.

%---------------------------------------------------------------------------
\subsection{Surface-aligned variants}
\label{sec:2dgs}

Because 3D Gaussians are volumetric, their zero-crossing does not define a
surface, and depth composited through a thick primitive is biased. 2D Gaussian
Splatting \cite{huang20242dgs} replaces each ellipsoid with an oriented planar
disk, so the primitive \emph{is} a surface element with a well-defined normal;
combined with normal-consistency and distortion regularisers it yields markedly
better geometry, and SuGaR \cite{guedon2024sugar} similarly binds Gaussians to an
extractable mesh.

Either representation supports mesh extraction by fusing rendered depth maps into
a truncated signed-distance volume \cite{curless1996tsdf} and running marching
cubes, which is attractive here because the object of interest is a relief surface
rather than a radiance field. Note that this route does not strictly require
surface-aligned primitives: composited depth from a 3D Gaussian model can be fused
in the same way, at the cost of the depth bias induced by volumetric primitives.

We evaluated 2DGS and report a negative result: on this capture it produced flat,
featureless renderings at PSNR~$\approx13$. The reason is instructive rather than
incidental. Surface-aligned primitives must recover an orientation as well as a
position, and disk normals are constrained mainly by observing a surface from
multiple directions. A single-sided capture supplies almost no such signal, so the
added constraint removes degrees of freedom the photometric loss was using without
supplying the multi-view evidence needed to determine the ones that remain. The
same regularisation that improves well-covered scenes is therefore harmful here.

%---------------------------------------------------------------------------
\subsection{Per-image appearance modelling}
\label{sec:appearance}

Coverage is the binding constraint (\S\ref{sec:sfm}), and the cheapest source of
additional coverage is additional clips of the same object. This is obstructed by
photometry rather than geometry: a handheld camera re-runs auto-exposure and
auto-white-balance per clip, so the same surface is recorded with different
responses. The model in \eqref{eq:compositing} has no parameter expressing ``same
radiance, different camera response'', so the photometric loss
\eqref{eq:photo} cannot be driven to zero anywhere; residual gradients stay high
across the whole image, and \S\ref{sec:densification} interprets uniformly high
$\tau_i$ as under-reconstruction and densifies without bound. Our three-clip merge
duly plateaued near PSNR~$14$ while a single clip reached~$24$.

NeRF-in-the-Wild \cite{martinbrualla2021nerfw} solves this by attaching a latent
appearance code $\bm{\ell}_k$ to each image $k$ and conditioning the emitted
radiance on it. The codes are free parameters optimised jointly with the scene ---
generative latent optimisation \cite{bojanowski2018glo} --- rather than the output
of an encoder. We use the same idea with a deliberately low-capacity decoder: a
small MLP maps $\bm{\ell}_k$ to a global affine colour transform applied to the
rendered image,
\begin{equation}
  \hat{\bm{I}}_k^{\text{corr}}(\bm{p}) \;=\; \bm{M}(\bm{\ell}_k)\,\hat{\bm{I}}_k(\bm{p}) \;+\; \bm{b}(\bm{\ell}_k),
  \qquad
  \bm{M}\in\mathbb{R}^{3\times3},\;\bm{b}\in\mathbb{R}^{3},
  \label{eq:appearance}
\end{equation}
with the decoder initialised so that $\bm{M}=\bm{I}$ and $\bm{b}=\bm{0}$;
training therefore starts exactly at the identity and departs from it only if
doing so reduces the loss.

\paragraph{Why this cannot absorb geometry.}
Capacity control is the crux: a per-image correction expressive enough to repaint
pixels individually would simply memorise each training view, destroying the
multi-view constraint that makes reconstruction possible. The transform
\eqref{eq:appearance} is applied with the \emph{same} $(\bm{M},\bm{b})$ at every
pixel, so for any spatial permutation $\pi$ of the pixels,
\begin{equation}
  \bm{M}\,\hat{\bm{I}}\big(\pi(\bm{p})\big)+\bm{b}
  \;=\;
  \Big(\bm{M}\,\hat{\bm{I}}(\bm{p})+\bm{b}\Big)\circ\pi ,
\end{equation}
i.e.\ the correction commutes with any rearrangement of image content. A map that
is invariant to where things are cannot encode where things are: it carries
exactly zero spatial information and can represent only a global photometric
change. It can absorb exposure, gain and white-balance drift, and nothing else.

At evaluation time we apply the transform decoded from the \emph{mean} training
latent, giving one canonical appearance. Fitting a held-out image's own latent
would leak information from the test image into the prediction and inflate the
reported metrics.

%---------------------------------------------------------------------------
\subsection{Pose refinement}
\label{sec:pose}

A residual error source is the poses themselves: \S\ref{sec:sfm} treats them as
fixed, yet low-overlap SfM yields noisier extrinsics, and a pose error is
indistinguishable from a geometry error to the photometric loss. Joint
optimisation of poses with the scene, as in BARF \cite{lin2021barf}, addresses
this, with the caveat that a high-frequency representation admits many poor local
minima --- BARF's central observation is that the optimisation must be coarse-to-fine.
The analogous mechanism in our setting is the progressive introduction of
spherical-harmonic bands (\S\ref{sec:3dgs}).

%---------------------------------------------------------------------------
\subsection{Evaluation}
\label{sec:eval}

We report PSNR, SSIM \cite{wang2004ssim} and LPIPS \cite{zhang2018lpips} on
held-out views. The three are deliberately complementary: PSNR derives from mean
squared error and is dominated by low-frequency agreement, rewarding a blurred
prediction; SSIM compares local luminance, contrast and structure statistics and
so is sensitive to the relief we care about; LPIPS compares deep features and
correlates better with human judgement of perceptual similarity.

Two methodological cautions apply to the comparisons in this report. First, model
selection must use validation views only --- selecting on the test split silently
converts it into a validation split. Second, these metrics are \emph{not}
comparable across reconstructions built at different image resolutions or with
different held-out splits: PSNR in particular shifts with the spatial frequency
content of the images. Where we compare a higher-resolution reconstruction against
a lower-resolution one we therefore compare \emph{distributions} over views ---
medians and the frequency of catastrophic views --- and state the caveat
explicitly, rather than reporting a single number as if it were a like-for-like
improvement.

Finally, a scalar mean conceals the failure structure that matters. Per-view
analysis of our reconstruction showed the errors were not uniformly distributed
but concentrated entirely in extreme close-ups at grazing incidence, which
identified a spatial-frequency deficit --- the source images had been downscaled
--- and directly motivated the higher-resolution reconstruction. Aggregate metrics
diagnose nothing; the distribution, and the images themselves, do.
